\documentclass[exercice,a4]{thedocument}%
\usepackage[french]{babel}%
\usepackage{pstricks}%
\usepackage{pstricks-add}%
\usepackage{tabularx}%
\usepackage{array}%
\startdatakeys
    \source{21GENMATAA2}
    \BO{2026}
    \cycle{4}
    \competencesDuSocle{CA,RA}
    \objectifsApprentissage{C4C1ac,C4C1ab}
    \motsCles{pourcentage,moyenne,tableur,lecture graphique}
\enddatakeys
\begin{document}%

\medskip

Une station de ski compte $30$ pistes. Ces pistes de ski sont soit vertes, soit bleues, soit rouges, soit noires. La couleur de la piste définit son niveau de difficulté pour skier.

\smallskip

Chaque piste de ski peut être soit ouverte, soit fermée.

\smallskip

Sur le site internet de la station de ski, on a pu trouver les informations suivantes :

\begin{center}
\begin{tabular}{|m{4cm} m{2.7cm}|m{4cm} m{2.7cm}|}\hline
\multicolumn{4}{|m{13.4cm}|}{\textbf{Lundi 17 février 2020}}\\ \hline
\multicolumn{2}{|m{6.7cm}|}{\textbf{Pistes vertes}}&\multicolumn{2}{|m{6.7cm}|}{\textbf{Pistes bleues}}\\
Les 7 pistes vertes sont toutes ouvertes.&\psset{unit=1cm}
\begin{pspicture}(-1.3,-1.3)(1.3,1.3)
\pscircle[fillstyle=solid,fillcolor=lightgray](0,0){1.3}
\psline(1.3;90)
\end{pspicture}&Nombre de pistes bleues : 8&\begin{pspicture}(-1.3,-1.3)(1.3,1.3)
\psline(1.3;90)
\pscircle(0,0){1.3}
\pswedge[fillstyle=solid,fillstyle=hlines](0,0){1.3}{90}{180}
\pswedge[fillstyle=solid,fillcolor=lightgray](0,0){1.3}{-180}{90}
\psframe(0,0)(-0.2,0.2)
\end{pspicture}\\ \hline
\multicolumn{2}{|c|}{\textbf{Pistes rouges}}&\multicolumn{2}{|c|}{\textbf{Pistes noires}}\\
Parmi les 10 pistes rouges, 8 pistes rouges sont ouvertes.&\begin{pspicture}(-1.3,-1.3)(1.3,1.3)
\pswedge[fillstyle=solid,fillcolor=lightgray](0,0){1.3}{162}{90}
\pswedge[fillstyle=hlines](0,0){1.3}{90}{162}
\end{pspicture}&Parmi les 5 pistes noires, 3 pistes noires sont ouvertes.&
\begin{pspicture}(-1.3,-1.3)(1.3,1.3)
\pswedge[fillstyle=solid,fillcolor=lightgray](0,0){1.3}{-126}{90}
\pswedge[fillstyle=hlines](0,0){1.3}{90}{234}
\end{pspicture}\\ \hline
\multicolumn{4}{|l|}{\emph{Information} : \psframe[fillstyle=hlines](1.2,0.5) \qquad\qquad Pistes fermées \qquad\qquad\qquad  \psframe[fillstyle=solid,fillcolor=lightgray](1.2,0.5) \qquad\qquad Pistes ouvertes}\rule[-3mm]{0mm}{12mm}\\ \hline
\end{tabular}
\end{center}

\medskip

\begin{enumerate}
\item Déterminer le nombre de pistes rouges fermées le lundi 17 février 2020.
\item Justifier qu’il y a six pistes bleues ouvertes le lundi 17 février 2020.
\item Parmi les pistes noires, quel est le pourcentage de pistes noires ouvertes le lundi 17 février 2020 ?
\item Le mercredi 19 février 2020, la nouvelle répartition affichée sur le site internet est la suivante :

\begin{center}
\begin{tabularx}{\linewidth}{|*{4}{>{\centering \arraybackslash}X|}}\hline
\textbf{Pistes vertes}&\textbf{Pistes bleues}&\textbf{Pistes rouges}&\textbf{Pistes noires}\\
Nombre de pistes : 7&Nombre de pistes : 8&Nombre de pistes : 10&Nombre de pistes : 5\\
Nombre de pistes ouvertes : 5&Nombre de pistes ouvertes : 4&Nombre de pistes ouvertes : 3&Nombre de pistes ouvertes : 1\\ \hline
\end{tabularx}
\end{center}

\medskip

Sur le site de la station on peut lire :

\og Votre forfait du jour est remboursé si plus de 50\,\% des pistes de la station sont fermées. \fg

Une cliente demande le remboursement de son forfait du jour du mercredi 19 février 2020.

La station de ski doit-elle effectuer ce remboursement ?
\item On a mesuré les hauteurs maximales de neige dans la station, exprimées en centimètre, pour chaque mois, de novembre 2018 à mars 2019.

On saisit ces mesures dans une feuille de calcul dont voici une copie d’écran :

\begin{center}
\begin{tabularx}{\linewidth}{|c|l|*{6}{>{\centering \arraybackslash}X|}}\hline
	&A				&B			&C		&D		&E		&F		&G\\ \hline
1	&				&Nov.	&Déc.&Janvier&Février&Mars	&\small Moyenne\\ \hline
2	&Saison 2018-2019&90		&120	&130	&120	&75		&107\\ \hline
3	&Saison 2019-2020&105		&130	&115	&140	&60		&\\ \hline
\end{tabularx}
\end{center}

	\begin{enumerate}
		\item Quelle formule a pu être saisie dans la cellule G2 avant d’être étirée jusqu’à la cellule G3 ?
		\item La moyenne des cinq hauteurs maximales de neige de la saison 2019-2020 est-elle supérieure à celle de la saison 2018-2019 ?
	\end{enumerate}
\end{enumerate}

\bigskip

\showthesource
\end{document}
